\begin{abstract}
A large literature predicts short-term price moves from the limit order book (LOB): signals such as order-flow
imbalance and the micro-price, queue and first-passage models, Hawkes and neural point processes, and deep networks
from convolutional models to Transformers and pretrained message-level encoders. Most of it stops at the forecast.
Models are scored on how often they guess the direction of the next move, and in an independent re-test of fifteen
published deep models on later Nasdaq data, F1 scores were between 48\% and 61\%, lower than those reported on FI-2010. A trader needs to know something else: whether the
information a model extracts survives the queue, the fill and the latency of a real order, and still makes money.

This survey follows a model from theory to production. It is three surveys in one, in the order the work is done:
order-book models---how each is built, what it predicts and which simple baselines it must beat; simulators---how a
model becomes a trading strategy and is tested against recorded markets with its orders in the real queues; and
hardware---how a tested strategy is deployed fast enough to run. Its thesis is that what matters for trading is the
expected profit of an action given the market state and the order's place in the queue, not the accuracy of the
forecast behind it.

The survey shows a practitioner how to check each step. The first check is whether a reported accuracy is real. In a
re-test of DeepLOB on a year of Apple data, with the same network, the construction of the label alone moved accuracy
from 65.9\% to 54.6\%; with the network held fixed, the choice of input decided most of the result in one large
comparison. The second check is whether a forecast survives execution: what a simulator must capture, which simulators
provide it, and a model-free execution benchmark that sets the floor. On a year of E-mini S\&P~500 data, resting
(leaving a limit order in the queue) and crossing (trading at once at the opposite price) cost similar amounts, about
$0.095$ ticks per contract, and cost rose with every step of delay. The third check is whether a model can be
deployed. Published inference times of order-book models range from 25\,\textmu s to 23\,ms, while a stage on the hot
path of a CME feed (the chain every market-data message passes through) must handle each packet in about
7.5\,\textmu s. A slower stage lets the queue in front of it grow: the strategy misses signals and its orders are
picked off (filled by faster traders after the price has moved). The survey sets out which stages FPGAs, GPUs and
other processors can accelerate.

The worked example throughout is \kasparhft{}\footnote{\kasparhft{} is available at
\url{https://github.com/vincent212/kaspar-hft}.}, an open-source, venue-independent trading system and
deterministic, queue-exact simulator, with handlers implemented for CME. Simulated orders enter the reconstructed first-in-first-out queues and are filled only by
recorded executions that reach them; a model is added as one more actor that sends its predictions to a strategy; and
the strategy tested in the simulator is the strategy deployed, with no change of code---moving from replay to paper and
live trading is a configuration setting.

Three principles run through the survey: the market is observed, not generated; complexity, in the model or in the
hardware, must earn its place by out-of-sample economic value; and representations are compared before
architectures.

\medskip
\noindent\textbf{Keywords:} limit order book, market microstructure, order-flow imbalance, micro-price,
Hawkes process, queue position, first-passage time, adverse selection, deterministic simulation, deep
learning, spatial-temporal Transformer, foundation models, uncertainty quantification, multiple testing,
Occam's razor, market simulation, reinforcement learning, hardware acceleration, FPGA, Kaspar-HFT.
\end{abstract}
